% GENERATED FILE. Edit canonical lessons, then run npm run generate:books.
% canonical-source-hash: fnv1a64:f39f1ea067eabaa3
% canonical-lessons: ES-C282-bienvenido, ES-C282-saludo, ES-C282-abrazo, ES-C282-sonrisa, ES-C282-invitado

\chapter{Meeting Someone --- Welcome, Greeting, Hug, Smile, Guest}
\label{ch:encuentro}
\hlchaptermodality{282}

\begin{chapteropening}
\textbf{By the end of this chapter:} \emph{I can welcome someone in Spanish and name the greeting, the hug, the smile and the guest that go with it.}

Complete ES-C282-invitado: say all five words in order, and say which two of them are participles.
\end{chapteropening}

\section[bienvenido]{bienvenido — well come, in two words you already have}

\label{lesson:ES-C282-bienvenido}



\begin{quote}
What did \emph{sucidus} describe before it meant dirty? (Unwashed wool, still greasy.) Now a word with nothing new in it but the joining.
\end{quote}

\subsection*{What to know first}
\begin{itemize}
\raggedright
  \item bien — well, and the first half of this word.
  \item hola — the greeting you have opened with since the start.
\end{itemize}

\begin{sounds}
\begin{itemize}
\raggedright
  \item \emph{bien-ve-NI-do}, stress on the \emph{ni}. The \emph{v} is the same soft sound you give \emph{b} — two letters, one sound. $\to$ reference
\end{itemize}
\end{sounds}

\begin{cousinweb}
\textbf{bienvenido} — \textquotedblleft{}welcome.\textquotedblright{}

Two pieces, both of them plain. \textbf{bien}, well, which you have; and \textbf{venido}, come, the participle of \emph{venir}. Latin put it the same way round: \emph{bene venutus}, well come.

English \textbf{welcome} looks like the identical build and is not one. Old English had \textbf{wilcuma}, from \emph{willa}, desire, and \emph{cuma}, a comer: the person you wanted to arrive. Norse and French pressure then bent \emph{wil-} into \emph{wel-}, and the word ended up resembling \textquotedblleft{}well\textquotedblright{} plus \textquotedblleft{}come\textquotedblright{} by accident.

So two languages reached the same shape from opposite directions. Spanish means it literally. English only appears to.
\end{cousinweb}

\begin{grammarlens}[title={What you've built}]
Nothing here was new except the joining.

\emph{Bien} you had already, and \emph{venido} has \emph{venir} audible inside it. A compound you can take apart is a word you do not have to memorise — and Spanish builds a great many of its politest words this way.
\end{grammarlens}

\subsection*{Your turn}
Say these aloud:

\begin{itemize}
\raggedright
  \item \emph{bienvenido}, then \emph{bienvenida}
  \item the two pieces on their own — \emph{bien}, \emph{venido}
  \item \emph{hola}, then \emph{bienvenido}, as you would at a door
\end{itemize}

\begin{tcolorbox}[breakable,colback=teal!4,colframe=teal!35!black,title={Before you move on}]
What are the two pieces of \emph{bienvenido}? (\emph{Bien} and \emph{venido}.) Is English \emph{welcome} built the same way? (No — it was \emph{wilcuma}, a desired comer.)
\end{tcolorbox}
\section[el saludo]{el saludo — a wish for your health, worn smooth}

\label{lesson:ES-C282-saludo}



\begin{quote}
Before the new one: what are the two pieces of \emph{bienvenido}? (\emph{Bien} and \emph{venido}.)
\end{quote}

\subsection*{What to know first}
\begin{itemize}
\raggedright
  \item bienvenido — what you say; this is what you are doing.
  \item hola — one particular instance of it.
\end{itemize}

\begin{sounds}
\begin{itemize}
\raggedright
  \item \emph{sa-LU-do}, stress on the \emph{lu}. The \emph{d} between two vowels is soft, closer to the \emph{th} of English \emph{this} than to a hard \emph{d}. $\to$ reference
\end{itemize}
\end{sounds}

\begin{cousinweb}
\textbf{el saludo} — \textquotedblleft{}a greeting.\textquotedblright{}

From \textbf{saludar}, to greet, and behind that Latin \textbf{salutem}, health — the accusative of \emph{salus}, soundness and safety. To greet a Roman was to wish him his health, and Spanish still raises a glass with \emph{salud} and answers a sneeze with it.

English took the same Latin word more than once. \textbf{Salute} is the soldier's version, borrowed whole. \textbf{Salutary} and \textbf{salubrious} kept the health and dropped the greeting altogether.

So a \emph{saludo} is not really a word. It is a wish, repeated until nobody hears it any more — which is what centuries do to every polite formula, in every language.
\end{cousinweb}

\begin{grammarlens}[title={What you've built}]
You have been greeting people since your first lesson without a name for the act.

This is the name. Having one lets you talk about the thing rather than only do it — and that is a step up in what you can say.
\end{grammarlens}

\subsection*{Your turn}
Say these aloud:

\begin{itemize}
\raggedright
  \item \emph{el saludo}
  \item \emph{un saludo}, then \emph{el saludo}
  \item \emph{bienvenido} and then \emph{el saludo} — the word, then the act
\end{itemize}

\begin{tcolorbox}[breakable,colback=teal!4,colframe=teal!35!black,title={Before you move on}]
What is behind \emph{saludo}? (\emph{Salutem}, health.) Which English word is the soldier's version of it? (\emph{Salute}.)
\end{tcolorbox}
\section[el abrazo]{el abrazo — the arm is still inside the word}

\label{lesson:ES-C282-abrazo}



\begin{quote}
Before the new one: what did \emph{saludo} wish you? (Your health.)
\end{quote}

\subsection*{What to know first}
\begin{itemize}
\raggedright
  \item la mano — a part of the body you already have.
  \item el saludo — the greeting this one belongs to.
\end{itemize}

\begin{sounds}
\begin{itemize}
\raggedright
  \item \emph{a-BRA-zo}, stress in the middle. The \emph{r} is a single tap, and the \emph{z} is the \emph{th} of \emph{think} in much of Spain and an \emph{s} elsewhere. $\to$ reference
\end{itemize}
\end{sounds}

\begin{cousinweb}
\textbf{el abrazo} — \textquotedblleft{}a hug.\textquotedblright{}

From \textbf{abrazar}, and inside it Latin \textbf{bracchium}, the arm: \emph{ad-bracchiare}, to put your arms around. The word says exactly what the body does, and Spanish never hid it — \emph{brazo} is still the word for the arm itself.

English \textbf{embrace} is the same build through French, \emph{in} plus \emph{bracchium}. A \textbf{bracelet} is a little arm-band. A \textbf{brace} holds a thing the way an arm holds it.

Latin took \emph{bracchium} from Greek \emph{brakhion}, which meant the shorter upper arm — from \emph{brakhys}, short, the same \emph{brachy-} that turns up in a scientist's vocabulary. The hug goes back to a word for shortness.
\end{cousinweb}

\begin{grammarlens}[title={What you've built}]
Three words now, and each of them is something you could point at.

A welcome, a greeting, an embrace. None of the three needed grammar to be useful — which is worth remembering when a language starts to feel like rules.
\end{grammarlens}

\subsection*{Your turn}
Say these aloud:

\begin{itemize}
\raggedright
  \item \emph{el abrazo}
  \item \emph{un abrazo}, which is how a letter ends
  \item \emph{el saludo} and \emph{el abrazo} one after the other
\end{itemize}

\begin{tcolorbox}[breakable,colback=teal!4,colframe=teal!35!black,title={Before you move on}]
What is inside \emph{abrazo}? (\emph{Bracchium}, the arm.) Which English word is the same build? (\emph{Embrace}.)
\end{tcolorbox}
\section[la sonrisa]{la sonrisa — a laugh with the volume down}

\label{lesson:ES-C282-sonrisa}



\begin{quote}
Before the new one: what is inside \emph{abrazo}? (\emph{Bracchium}, the arm.)
\end{quote}

\subsection*{What to know first}
\begin{itemize}
\raggedright
  \item el abrazo — the last one, and another thing a face or a body does.
  \item bien — well, which you have had since the beginning.
\end{itemize}

\begin{sounds}
\begin{itemize}
\raggedright
  \item \emph{son-RI-sa}, stress on the \emph{ri}. The \emph{s} is clear and never buzzes into a \emph{z} sound the way an English \emph{s} between vowels does. $\to$ reference
\end{itemize}
\end{sounds}

\begin{cousinweb}
\textbf{la sonrisa} — \textquotedblleft{}a smile.\textquotedblright{}

From \textbf{sonreír}, and that is Latin \textbf{subridere}: \emph{sub-}, under or slightly, plus \textbf{ridere}, to laugh. A smile is a laugh with the volume turned down. The \emph{sub-} wore to \emph{son-} on the way in.

English keeps \emph{ridere} only at the loud end, and only in loans: to \textbf{deride} is to laugh at, \textbf{ridicule} is the laughing itself, and something \textbf{risible} is worth a laugh.

So Spanish built its quiet word out of the noisy one by putting \emph{sub-} in front. English never built it at all — \emph{smile} is Germanic and has no Latin in it.
\end{cousinweb}

\begin{grammarlens}[title={What you've built}]
\emph{Sub-} is a dial, not a word.

Under, slightly, a little less than the whole thing. Once you can hear it, a long Latin word stops being one item to memorise and becomes two smaller ones you can already handle.
\end{grammarlens}

\subsection*{Your turn}
Say these aloud:

\begin{itemize}
\raggedright
  \item \emph{la sonrisa}
  \item \emph{una sonrisa}, then \emph{la sonrisa}
  \item \emph{el abrazo} and \emph{la sonrisa}, one after the other
\end{itemize}

\begin{tcolorbox}[breakable,colback=teal!4,colframe=teal!35!black,title={Before you move on}]
What does \emph{sub-} do to \emph{ridere}? (Turns the laugh down to a smile.) Which English words kept the laughing? (\emph{Deride}, \emph{ridicule}, \emph{risible}.)
\end{tcolorbox}
\section[el invitado]{el invitado — named by what was done to him}

\label{lesson:ES-C282-invitado}



\begin{quote}
Before the last of the five: what does \emph{sub-} do to \emph{ridere}? (Turns the laugh down to a smile.) And what is inside \emph{abrazo}? (The arm.)
\end{quote}

\subsection*{What to know first}
\begin{itemize}
\raggedright
  \item el abrazo — what you might give this person.
  \item la sonrisa — and what you would give with it.
\end{itemize}

\begin{sounds}
\begin{itemize}
\raggedright
  \item \emph{in-vi-TA-do}, stress on the \emph{ta}. Four syllables, all of them short and even; Spanish does not swallow the unstressed ones the way English does. $\to$ reference
\end{itemize}
\end{sounds}

\begin{cousinweb}
\textbf{el invitado} — \textquotedblleft{}a guest.\textquotedblright{}

From \textbf{invitar}, Latin \textbf{invitare}, to invite. \emph{Invitado} is its participle, \textquotedblleft{}invited\textquotedblright{} — so Spanish names the guest by what was done to him, exactly as \emph{bienvenido} names a welcome by what is wished.

Put the two side by side. \textbf{venido} in the first of these five and \textbf{invitado} in the last are the same piece of grammar doing the same job: an action turned into a description of a person.

English \textbf{guest} is a different word altogether, Germanic \emph{gastiz} — which is also the source of Latin \emph{hostis}, a stranger and then an enemy. Two peoples had one word for the man at the door and could not agree whether he was a visitor or a threat.
\end{cousinweb}

\begin{grammarlens}[title={What you've built}]
Five, and the run is closed: \emph{bienvenido}, \emph{el saludo}, \emph{el abrazo}, \emph{la sonrisa}, \emph{el invitado}.

Two of the five are participles, and that is the pattern worth carrying forward. Spanish is unusually willing to name a person or a thing by an action already finished — the one invited, the one welcomed — and once you notice it you will find it everywhere.
\end{grammarlens}

\subsection*{Your turn}
Say these aloud:

\begin{itemize}
\raggedright
  \item \emph{el invitado}, then \emph{la invitada}
  \item all five in order — \emph{bienvenido}, \emph{el saludo}, \emph{el abrazo}, \emph{la sonrisa}, \emph{el invitado}
  \item the two participles on their own — \emph{bienvenido}, \emph{invitado}
\end{itemize}

\begin{tcolorbox}[breakable,colback=teal!4,colframe=teal!35!black,title={Before you move on}]
What is an \emph{invitado}, literally? (One who has been invited.) Which other word in this run is a participle? (\emph{Bienvenido} — well come.) And what did \emph{saludo} wish you? (Your health.)
\end{tcolorbox}
